\chapter{Introduction}

\section{Contexte et objectif}
Cet atelier s'inscrit dans le module Java EE du cycle ingénieur LSI de la FST de Tanger, sous la direction de
\encadrant. Le sujet demande de construire une application web de commerce électronique en respectant deux
contraintes d'architecture : le patron \textbf{MVC2}, où les servlets jouent le rôle de contrôleurs et les pages
JSP celui de vues, et la persistance par \textbf{JPA}, ici avec EclipseLink, au lieu de requêtes JDBC écrites à la main
comme à l'atelier~1.

L'application, baptisée \textit{Atlas Boutique}, reprend l'identité visuelle de l'atelier~1 (barre latérale sombre,
vert profond, safran pour l'élément courant) avec une interface plus soignée. Elle fonctionne sans aucune
dépendance externe au moment de l'exécution : la bibliothèque jQuery est livrée avec l'application, il n'y a ni
CDN ni police distante.

\section{Périmètre retenu}
Le sujet cite un grand nombre de classes. Il recommande lui-même de se concentrer sur le panier, la vitrine et
l'internaute. J'ai donc retenu la « version A » du périmètre :

\begin{itemize}
  \item l'\textbf{internaute}, spécialisé en \textbf{client} et \textbf{administrateur} ;
  \item la \textbf{vitrine} et ses \textbf{catégories} ;
  \item les \textbf{produits} ;
  \item le \textbf{panier} et ses \textbf{lignes}.
\end{itemize}

La \textbf{commande} n'est pas réalisée. Ce choix est assumé : une commande bien faite exige de décrémenter le
stock dans une transaction avec contrôle de concurrence, ce qui mérite un atelier à part entière. Le
chapitre~\ref{chap:perspectives} décrit comment elle s'ajouterait sans remettre en cause l'existant.

\section{Démarche}
Trois principes ont guidé le travail.
\begin{enumerate}
  \item \textbf{Décider, puis justifier.} Chaque choix structurant (stratégie d'héritage, verrouillage, gestion des
        transactions) est expliqué par son pourquoi et comparé à au moins une alternative.
  \item \textbf{Vérifier à chaque étape.} Chaque couche a été validée avant de passer à la suivante : tables et
        contraintes dans MySQL, contrôles automatiques, puis parcours dans le navigateur.
  \item \textbf{Sécurité par défaut.} Les protections (jeton CSRF, échappement, requêtes paramétrées, en-têtes) font
        partie de la structure de l'application, elles ne sont pas ajoutées à la fin.
\end{enumerate}

\section{Taille du projet}
Le projet compte environ 4\,000 lignes de Java pour le code de production (45 classes), 500 lignes pour les
tests (7 classes), 830 lignes pour les vues et les gabarits JSP, et 580 lignes de CSS et de JavaScript.

\section{Organisation du rapport}
Le chapitre~\ref{chap:conception} présente la conception (classes, règles métier, choix d'héritage). Le
chapitre~\ref{chap:env} décrit l'environnement et le projet Maven. Les chapitres~\ref{chap:persistance}
et~\ref{chap:services} traitent de la persistance et des services. Les chapitres~\ref{chap:web} et~\ref{chap:vues}
détaillent la couche web et les vues. La sécurité fait l'objet du chapitre~\ref{chap:securite}, les tests du
chapitre~\ref{chap:tests} et la gestion de version du chapitre~\ref{chap:git}. Le rapport se termine par les
limites, les perspectives et une conclusion.
